\documentclass[10pt,a4paper]{amsart}
\usepackage[margin=19mm]{geometry}
\usepackage{amsmath,amssymb,mathtools}
\usepackage[scr=rsfs,cal=euler]{mathalfa}
\usepackage{xfrac}
\usepackage[unicode,hidelinks,bookmarks=false]{hyperref}
\newcommand{\ie}{i.e.\ }
\newcommand{\cf}{cf.\ }
\newcommand{\resp}{respectively\ }
\newcommand{\Subsection}{\subsection}
\providecommand{\nobreakdash}{\nobreak\hbox{-}}
\setlength{\emergencystretch}{2em}
\multlinegap0pt
\usepackage{bm}
\usepackage{bbm}
\usepackage{dsfont}
\usepackage{xspace}
\usepackage{cancel}
\usepackage{mathtools}
\usepackage{amsthm}
\makeatletter
\@ifundefined{theorem}{\newtheorem{theorem}{Theorem}}{}
\makeatother
\DeclareMathOperator {\myRD}{RD}
\title[On the arithmetic properties of partitions into parts simult]{On the arithmetic properties of partitions into parts simultaneously $4$-regular and $9$-distinct}
\author{Mohammed L. Nadji, Moussa Ahmia}
\date{Source: arXiv:2506.07704v2 (CC BY 4.0)}
\begin{document}
\maketitle

\noindent\emph{Source and licence.} On the arithmetic properties of partitions into parts simultaneously $4$-regular and $9$-distinct. Version arXiv:2506.07704v2 (CC BY 4.0), distributed under the Creative Commons Attribution 4.0 International licence, as recorded in the arXiv licence metadata for this exact version and shown on the abstract page https://arxiv.org/abs/2506.07704v2. Full rights notice: licenses/arxiv-2506.07704-CC-BY-4.0.txt.\par\smallskip

\noindent\textit{Editorial scope.} Counted unit: the Theorem whose source label is \texttt{thm} in the article (source0.tex lines 320--325, source label \texttt{thm}), delivered together with the article's own complete original proof of it (source0.tex lines 327--369, one \texttt{proof} environment of 43 lines). No mathematics is written, repaired or re-derived: statement and proof are the article's own text line for line. Required context retained uncounted: psidissection (lines 167--170); binomiallemma (lines 192--194); d2 (lines 207--212). Every definition, display and label of those lines is retained unchanged. Omitted from this package and not redistributed: 1--166, 171--191, 195--206, 213--319, 326--326, 370--476. Nothing inside a retained excerpt is omitted or altered: every retained body is delivered exactly as the article's lines are written, so no omission receipt of any kind accompanies this package. Citations: the retained excerpts carry no citation command at all, so no bibliography entry of the article is transcribed and no reference is redistributed. Reference adaptation: none was needed and none was made; every reference inside the delivered unit resolves inside it, so no unresolved reference or citation is produced. Printed slips kept literal: no printed word, symbol, hypothesis or number of the counted statement or of its proof was altered. Layout and class: the article's own class is \texttt{amsart} with packages including hyperref, amsmath, babel, inputenc, amssymb, amsthm, graphics, graphicx, enumerate, array, bm, a4wide, color, url; the wrapper is the lane's pinned \texttt{amsart} editorial wrapper at 10pt on A4, which restates the theorem-like environments the retained text needs and adds the article's own macro definitions that the retained text needs (\texttt{\textbackslash myRD}), with extra wrapper packages (bm, bbm, dsfont, xspace, cancel, mathtools). The delivered numbering is the unit's own. Assets: no retained span contains a figure environment, a graphics-inclusion command, a PDF-page inclusion, a table or a TikZ picture; no image file is copied into this package, the package declares no asset dependency and no image file is redistributed with it. Licence: version arXiv:2506.07704v2 of this article is distributed under the Creative Commons Attribution 4.0 International licence, as recorded in the arXiv licence metadata for this exact version and shown on the abstract page https://arxiv.org/abs/2506.07704v2. A full rights notice is delivered at licenses/arxiv-2506.07704-CC-BY-4.0.txt. Characters: every retained excerpt is pure ASCII, so the delivered engine, XeLaTeX with its default Latin Modern text font, reports no missing character.
\begin{equation}
\psi(q)= \sum_{k=0}^{(p-3)/2} q^{k(k+1)/2} f(q^{\frac{p^{2}+(2k+1)p}{2}}, q^{\frac{p^{2}-(2k+1)p}{2}})+ q^{\frac{p^{2}-1}{8}}\psi(q^{p^{2}}).
 \label{psidissection}
\end{equation}

\begin{equation}
 f_{pm}^{p^{k-1}} \equiv f_{m}^{p^{k}}  \ (\mathrm{mod} \ p^{k}). \label{binomiallemma} 
\end{equation}

\begin{align}
&\sum_{n\geq 0} \myRD^{(4, 9)}(2n)q^{n}=\dfrac{f_{6}^{7}f_{9}^{7}}{f_{3}^{9}f_{18}^{5}}+2q\dfrac{f_{6}^{6}f_{9}^{4}}{f_{3}^{8}f_{18}^{2}}+4q^{2}\dfrac{f_{6}^{5}f_{9}f_{18}}{f_{3}^{7}}, \label{dd1}\\
&\sum_{n\geq 0} \myRD^{(4, 9)}(6n+2)q^{n}=2\dfrac{f_{2}^{6}f_{3}^{4}}{f_{1}^{8}f_{6}^{2}}, \label{d1}\\
&\sum_{n\geq 0} \myRD^{(4, 9)}(6n+4)q^{n}=4\dfrac{f_{2}^{5}f_{3}f_{6}}{f_{1}^{7}}, \label{d2}\\
&\sum_{n\geq 0} \myRD^{(4, 9)}(6n+5)q^{n}= 6\dfrac{f_{2}^{2}f_{3}^{2}f_{6}^{2}}{f_{1}^{6}}. \label{d3}
\end{align}

 \begin{theorem} \label{thm} For any prime $p\equiv 5 \ (\mathrm{mod} \ 6)$, $\alpha \geq 0$, and $n\geq 0$, we have
\[
\myRD^{(4, 9)} \big(6p^{2\alpha+1}(pn+i)+3p^{2\alpha+2}+1 \big)\equiv 0 \ (\mathrm{mod} \ 12),
\]
for all $1\leq i \leq p-1$.
\end{theorem}

\begin{proof} 
In view of \eqref{binomiallemma} and \eqref{d2}, with $p=3$ and $k=1$, we find that
\begin{equation}
\sum_{n\geq 0} \myRD^{(4, 9)}(6n+4)q^{n}=4\dfrac{f_{2}^{5}f_{3}f_{6}}{f_{1}^{7}} \equiv 4 \dfrac{f_{2}^{2} f_{6}^{2}}{f_{1} f_{3}} = 4\psi(q) \psi(q^{3}) \ (\mathrm{mod} \ 12). \label{lel1}
\end{equation}
Define
\begin{equation} 
\sum_{n\geq 0} b(n)q^{n} = \psi(q) \psi(q^{3}). \label{e1}
\end{equation}
From  Equations \eqref{lel1} and \eqref{e1} we have 
\begin{equation}
\myRD^{(4, 9)}(6n+4) \equiv 4b(n) \ (\mathrm{mod} \ 12).  \label{e2}
\end{equation} 
Now, consider the congruence equation
\[
\frac{k^{2}+k}{2}+3\cdot \frac{m^{2}+m}{2} \equiv  \frac{4p^{2}-4}{8} \ (\mathrm{mod} \ p).
\]
which is equivalent to
\begin{equation}
(2k+1)^{2}+ 3\cdot (2m+1)^{2}  \equiv  0 \ (\mathrm{mod} \ p), \label{e3}
\end{equation}
where $0 \leq k,m \leq (p-1)/2$ and $p$ is a prime number such that ($\frac{-3}{p}$)$=-1$. Since ($\frac{-3}{p}$)$=-1$ for $p \equiv 5 \ (\mathrm{mod} \ 6)$, the congruence relation of Equation \eqref{e3} holds if and only if both $k=m=(p-1)/2$. Substitute Equation \eqref{psidissection} into \eqref{e3} and extract the terms in which the powers of $q$ are congruent to $(p^{2}-1)/2$ modulo $p$, and then divide by $q^{\frac{p^{2}-1}{2}}$, we find that
\[
\sum_{n\geq 0} b \left(pn+\frac{p^{2}-1}{2} \right) q^{pn} = \psi(q^{p^{2}}) \psi(q^{3p^{2}}),
\]
which implies that 
\[
\sum_{n\geq 0} b\left( p^{2}n+\frac{p^{2}-1}{2} \right)q^{n} = \psi(q) \psi(q^{3}),
\]
and for $n\geq 0$,
\begin{equation}
b \left(p^{2}n+pi+\frac{p^{2}-1}{2}\right) = 0, \label{e5}
\end{equation}
where  $1\leq i \leq p-1$. By induction, we obtain  that for all $n, \alpha \geq 0$,
\begin{equation}
b\left( p^{2\alpha}n+\frac{p^{2\alpha}-1}{2} \right)=b(n). \label{e6}
\end{equation}
Replacing $n$ by $p^{2}n+pi+\frac{p^{2}-1}{2}$ ($1\leq i \leq p-1$) in \eqref{e6} and using \eqref{e5}, we deduce that for $n\geq 0$ and $\alpha \geq 0$,
\[
b\left( p^{2\alpha+2}n+p^{2\alpha+1}i+\frac{p^{2\alpha+2}-1}{2} \right)=0.
\]
Replacing $n$ by $p^{2\alpha+2}n+p^{2\alpha+1}i+\frac{p^{2\alpha+2}-1}{2}$ in Equation \eqref{e2}, we obtain the desired result.
\end{proof}

\end{document}
